% ============================================================
%  DOCUMENTO TÉCNICO — CHATBOT FACULTAD
%  Arquitectura, código y operación del sistema RAG
%  Facultad de Inteligencia Artificial e Ingenierías
%  Universidad de Caldas
%  Compilar con: pdflatex documento-tecnico.tex  (×2)
% ============================================================

\documentclass[11pt, a4paper, oneside]{report}

\usepackage[utf8]{inputenc}
\usepackage[T1]{fontenc}
\usepackage[spanish,es-tabla]{babel}
\usepackage{lmodern}
\usepackage{microtype}

\usepackage[
  top=2.4cm, bottom=2.4cm,
  left=2.6cm, right=2.4cm,
  headheight=16pt
]{geometry}

\usepackage[table, dvipsnames]{xcolor}
\definecolor{navy}{HTML}{00407D}
\definecolor{navydark}{HTML}{002D59}
\definecolor{navysoft}{HTML}{EDF3F9}
\definecolor{naranja}{HTML}{F27022}
\definecolor{naranjasoft}{HTML}{FFF4EC}
\definecolor{grisclaro}{HTML}{F5F7FA}
\definecolor{gristexto}{HTML}{333333}
\definecolor{grismuted}{HTML}{667689}
\definecolor{codebg}{HTML}{F0F4F8}

\usepackage{helvet}
\renewcommand{\familydefault}{\sfdefault}

\usepackage{graphicx}
\usepackage{float}
\usepackage{booktabs}
\usepackage{tabularx}
\usepackage{array}
\usepackage{multirow}
\usepackage{listings}
\usepackage{tikz}
\usetikzlibrary{arrows.meta, positioning, fit, calc, shapes.multipart}
\usepackage[most]{tcolorbox}
\usepackage{fancyhdr}
\usepackage{enumitem}
\usepackage{titlesec}
\usepackage{tocbibind}
\usepackage[
  colorlinks=true,
  linkcolor=navy,
  urlcolor=naranja,
  citecolor=navy,
  pdftitle={Documento Técnico — Chatbot Facultad},
  pdfauthor={Facultad de Inteligencia Artificial e Ingenierías — Universidad de Caldas},
  pdfsubject={Arquitectura y documentación técnica del asistente virtual RAG}
]{hyperref}

\newcolumntype{L}[1]{>{\raggedright\arraybackslash}p{#1}}
\newcolumntype{C}[1]{>{\centering\arraybackslash}p{#1}}
\newcolumntype{Y}{>{\raggedright\arraybackslash}X}

\setcounter{tocdepth}{2}
\setlength{\parskip}{6pt}
\setlength{\parindent}{0pt}

% Listings
\lstset{
  basicstyle=\ttfamily\small\color{gristexto},
  backgroundcolor=\color{codebg},
  frame=single,
  rulecolor=\color{navy!30},
  breaklines=true,
  columns=fullflexible,
  keepspaces=true,
  showstringspaces=false,
  aboveskip=8pt,
  belowskip=8pt,
  xleftmargin=4pt,
  xrightmargin=4pt
}

\newtcolorbox{notabox}{
  colback=navysoft, colframe=navy,
  fonttitle=\bfseries\color{navy}, title={Nota técnica},
  arc=5pt, boxrule=0.8pt, left=8pt, right=8pt, top=5pt, bottom=5pt
}
\newtcolorbox{advertenciabox}{
  colback=naranjasoft, colframe=naranja,
  fonttitle=\bfseries\color{naranja}, title={Importante},
  arc=5pt, boxrule=0.8pt, left=8pt, right=8pt, top=5pt, bottom=5pt
}
\newtcolorbox{archibox}[1]{
  colback=grisclaro, colframe=navy,
  fonttitle=\bfseries\color{white}, title={#1},
  coltitle=white,
  attach boxed title to top left={yshift=-2mm, xshift=4mm},
  boxed title style={colback=naranja, arc=3pt},
  arc=5pt, left=8pt, right=8pt, top=10pt, bottom=6pt
}

\pagestyle{fancy}
\fancyhf{}
\fancyhead[L]{\small\textcolor{navy}{\nouppercase{\leftmark}}}
\fancyhead[R]{\small\textcolor{naranja}{\bfseries Documento técnico}}
\fancyfoot[C]{\small\textcolor{grismuted}{\thepage}}
\renewcommand{\headrulewidth}{1.2pt}
\renewcommand{\headrule}{%
  \hbox to\headwidth{%
    \color{navy}\leaders\hrule height 1.2pt\hfill
    \color{naranja}\vrule width 2.2cm height 1.2pt
  }%
}

\titleformat{\chapter}[display]
  {\normalfont}
  {\textcolor{naranja}{\rule{1.8cm}{3.5pt}}\\[0.35cm]
   {\large\bfseries\textcolor{navy}{Capítulo \thechapter}}}
  {0.35cm}
  {\Huge\bfseries\textcolor{navy}}
  [\vspace{0.25cm}{\color{naranja}\titlerule[2pt]}\vspace{0.15cm}{\color{navy!40}\titlerule[0.4pt]}]
\titlespacing*{\chapter}{0pt}{-12pt}{18pt}

\titleformat{\section}
  {\color{navy}\Large\bfseries}
  {\thesection}{0.55em}{}
  [\vspace{-2pt}{\color{naranja}\titlerule[1.2pt]}\vspace{4pt}]

\titleformat{\subsection}
  {\color{navy}\large\bfseries}
  {\thesubsection}{0.45em}{}

\newcommand{\ruta}[1]{\texttt{\textcolor{navy}{#1}}}
\newcommand{\captura}[2]{%
  \begin{figure}[H]
    \centering
    \begin{tcolorbox}[
      enhanced, colback=navysoft, colframe=navy, boxrule=1pt, arc=8pt,
      width=0.94\linewidth, height=5.4cm, valign=center, halign=center,
      shadow={1mm}{-1mm}{0mm}{black!12},
      borderline west={4pt}{0pt}{naranja}
    ]
      {\large\textcolor{navy}{\bfseries[ Espacio para captura / diagrama ]}}\\[0.45cm]
      {\normalsize\textcolor{gristexto}{#2}}\\[0.55cm]
      {\footnotesize\texttt{documentation/img/#1.png}}
    \end{tcolorbox}
    \caption{#2}
    \label{fig:#1}
  \end{figure}%
}

\begin{document}

% ───────────────────────── PORTADA ─────────────────────────
\begin{titlepage}
  \newgeometry{margin=0cm}
  \begin{tikzpicture}[remember picture, overlay]
    % Fondo navy
    \fill[navy] (current page.south west) rectangle (current page.north east);
    % Franja inferior en azul más claro
    \fill[navy!70!white] ([yshift=0cm]current page.south west)
      rectangle ([yshift=1.6cm]current page.south east);
    % Franja diagonal decorativa (azul claro)
    \fill[white, opacity=0.08]
      ([yshift=6cm]current page.south west) --
      ([yshift=9.5cm]current page.south west) --
      ([yshift=4cm]current page.south east) --
      ([yshift=0.5cm]current page.south east) -- cycle;
    % Barra lateral azul
    \fill[navydark] (current page.north west)
      rectangle ([xshift=0.55cm]current page.south west);
  \end{tikzpicture}
  
  \color{white}
  \centering
  
  % Logo sin fondo blanco (transparente)
  \vspace*{1.8cm}
  \includegraphics[height=2.2cm]{logo-facultad1.png}
  \vspace{1.0cm}
  
  {\Large \bfseries Facultad de Inteligencia Artificial e Ingenierías}\\[0.15cm]
  {\large Universidad de Caldas}\\[2.0cm]
  
  {\color{naranja}\rule{0.5\linewidth}{2pt}}\\[0.6cm]
  
  {\fontsize{30}{36}\selectfont\bfseries Documento Técnico}\\[0.5cm]
  
  {\LARGE \bfseries Arquitectura y código del
Asistente Virtual}\\[1cm]
  
  {\large Daniela Agudelo Martínez}\\[1cm]
  {\normalsize Manizales, Colombia}
  
  \vfill
  
  {\small Versión 1.0 \quad\textcolor{naranja}{|}\quad \today}\\[2.0cm]
  
  \restoregeometry
\end{titlepage}

\color{gristexto}
\pagecolor{white}

\newpage
\thispagestyle{empty}
\section*{Control de versiones}
\begin{tabularx}{\linewidth}{C{1.6cm} C{2.8cm} L{3.8cm} X}
  \toprule
  \rowcolor{navy}
  \textcolor{white}{\textbf{Versión}} &
  \textcolor{white}{\textbf{Fecha}} &
  \textcolor{white}{\textbf{Autor}} &
  \textcolor{white}{\textbf{Descripción}} \\
  \midrule
  1.0 & \today & Daniela Agudelo & Documento técnico inicial \\
  \bottomrule
\end{tabularx}

\tableofcontents
\newpage
\listoffigures

% ============================================================
\chapter{Introducción}

\section{Objetivo}
Este documento describe la \textbf{arquitectura}, los \textbf{módulos de código},
los \textbf{contratos de API}, la \textbf{seguridad} y el \textbf{despliegue}
del Asistente Virtual, para facilitar
mantenimiento, auditoría académica y transferencia tecnológica a la Facultad.

\section{Alcance}
\begin{itemize}
  \item Backend RAG.
  \item Panel de administración.
  \item Plugin WordPress.
\end{itemize}

\section{Stack tecnológico}
\begin{tabularx}{\linewidth}{L{3.5cm} L{3.8cm} Y}
  \toprule
  \rowcolor{navy}
  \textcolor{white}{\textbf{Capa}} &
  \textcolor{white}{\textbf{Tecnología}} &
  \textcolor{white}{\textbf{Rol}} \\
  \midrule
  API RAG & FastAPI, LlamaIndex & Orquestación, chat, admin \\
  \rowcolor{navysoft}
  Vectores & ChromaDB & Índice semántico persistente \\
  LLM & NVIDIA NIM / Gemini & Generación de respuestas \\
  \rowcolor{navysoft}
  Embeddings & NVIDIA API & Vectorización de documentos y consultas \\
  Admin UI & React, Vite, TS & Panel operativo \\
  \rowcolor{navysoft}
  Admin DB & PostgreSQL & Usuarios, settings, audit \\
  Sitio público & WordPress + React & Widget de chat \\
  \bottomrule
\end{tabularx}

% ============================================================
\chapter{Arquitectura del sistema}

\section{Vista lógica}
El sistema sigue un patrón de tres superficies:

\begin{enumerate}
  \item \textbf{Cliente público} $\rightarrow$ proxy REST de WordPress
        $\rightarrow$ \ruta{POST /chat}.
  \item \textbf{Panel admin} $\rightarrow$ JWT $\rightarrow$ rutas \ruta{/auth/*}
        y \ruta{/admin/*}.
  \item \textbf{Motor RAG} (servicios internos): carga, chunking, embeddings,
        recuperación, generación.
\end{enumerate}

\begin{figure}[H]
\centering
\begin{tikzpicture}[
  font=\small\sffamily,
  box/.style={draw=navy, fill=navysoft, thick, rounded corners=3pt,
              align=center, minimum width=3.1cm, minimum height=1.05cm},
  orbox/.style={draw=naranja, fill=naranjasoft, thick, rounded corners=3pt,
                align=center, minimum width=3.1cm, minimum height=1.05cm},
  arr/.style={-{Latex}, thick, color=navy}
]
  \node[box] (wp) {Widget\\WordPress};
  \node[box, right=1.1cm of wp] (proxy) {Proxy REST\\WP};
  \node[orbox, right=1.1cm of proxy] (api) {FastAPI\\API RAG};
  \node[box, below=1.1cm of api] (chroma) {ChromaDB};
  \node[box, left=1.1cm of chroma] (llm) {LLM\\NIM / Gemini};
  \node[orbox, above=1.1cm of api] (admin) {Panel Admin\\React};
  \node[box, right=1.1cm of api] (pg) {PostgreSQL\\Admin};

  \draw[arr] (wp) -- (proxy);
  \draw[arr] (proxy) -- (api);
  \draw[arr] (admin) -- (api);
  \draw[arr] (api) -- (chroma);
  \draw[arr] (api) -- (llm);
  \draw[arr] (api) -- (pg);
\end{tikzpicture}
\caption{Arquitectura lógica de alto nivel}
\label{fig:arq-logica}
\end{figure}

\section{Flujo RAG}
\begin{enumerate}
  \item El usuario envía un mensaje con {session\_id}.
  \item Se recupera o crea el motor conversacional en memoria.
  \item El retriever consulta Chroma con embeddings NVIDIA.
  \item Los fragmentos (\textit{top\_k}) se envían al LLM junto con el historial condensado.
  \item La API responde con answer y sources[] (sin citas embebidas en el texto).
\end{enumerate}

\begin{figure}[H]
  \centering
  \includegraphics[width=0.92\linewidth]{diagrama-flujo.png}
  \caption{Diagrama detallado del flujo de chat RAG}
  \label{fig:documentos-subida}
\end{figure}

\section{Flujo de indexación}
Solo administradores autenticados:

\begin{enumerate}
  \item Subida/eliminación de archivos en data/ vía /admin/documents/*.
  \item Se marca reindex\_required = true.
  \item POST /admin/reindex ejecuta IngestService.run():
        carga $\rightarrow$ parseo/chunking $\rightarrow$ \\embeddings $\rightarrow$ upsert Chroma.
  \item Se invalidan caché de índice, retriever compartido y sesiones de chat.
\end{enumerate}

% ============================================================
\chapter{Backend: chatbot-facultad}

\section{Estructura de carpetas}
\begin{lstlisting}
chatbot-facultad/
  api/                 # FastAPI: health, chat, routers admin
  admin/               # Auth JWT, settings, docs, users
  alembic/             # Migraciones PostgreSQL
  llm/                 # Adaptadores Gemini / NVIDIA
  services/            # ChatService, IngestService
  generation/          # Engine, prompt, session_store
  carga_documentos/    # Loader, parsers, pipeline
  retrieval/           # Retriever y fuentes
  storage/             # Chroma db
  tests/               # pytest
  data/                # Corpus documental
  chroma_db/           # Persistencia vectorial
\end{lstlisting}

\section{Punto de entrada}
\begin{archibox}{api/main.py}
Registra CORS, lifespan (carga de API keys desde BD + {configure\_settings}),
routers {/auth} y {/admin}, y endpoints públicos
{GET /health} y {POST /chat}.
\end{archibox}

\section{Módulo LLM (adaptadores)}
\begin{tabularx}{\linewidth}{L{5cm} Y}
  \toprule
  \rowcolor{navy}
  \textcolor{white}{\textbf{Archivo}} &
  \textcolor{white}{\textbf{Responsabilidad}} \\
  \midrule
  \ruta{llm/base.py} & Interfaz del proveedor. \\
  \rowcolor{navysoft}
  \ruta{llm/factory.py} & Selección {gemini\textbar nvidia} según entorno/settings. \\
  \ruta{llm/nvidia\_provider.py} & Configura LlamaIndex NVIDIA NIM. \\
  \rowcolor{navysoft}
  \ruta{llm/gemini\_provider.py} & Configura Google GenAI. \\
  \bottomrule
\end{tabularx}

\section{Ingesta y almacenamiento}
\begin{itemize}
  \item Extensiones:
  {.pdf}, {.txt}, {.md}.
  \item Chunking alineado a $\sim$512 tokens con solape configurable.
  \item Persistencia en colección Chroma {faculty\_docs}.
\end{itemize}

\section{Generación y sesiones}
\begin{archibox}{generation/}
{engine\_factory.py} construye el chat engine; {session\_store.py}
mantiene \\ engines/memorias en memoria con TTL ({SESSION\_TTL\_HOURS}).
{prompt.py} define el system prompt restrictivo (sin alucinaciones ni citas
en el texto).
\end{archibox}

\section{Configuración}
Centraliza rutas, defaults de chunking, top\_k, memoria y resolución de API keys:

\begin{enumerate}
  \item Clave en memoria cargada desde BD (cifrada/descifrada).
  \item Fallback a variables de entorno {GEMINI\_API\_KEY} / {NVIDIA\_API\_KEY}.
\end{enumerate}

% ============================================================
\chapter{Panel de administración (API)}

\section{Autenticación}
\begin{itemize}
  \item Login: POST /auth/login} $\rightarrow$ JWT Bearer.
  \item Caducidad: {JWT\_EXPIRE\_HOURS} (default 12).
  \item Contraseñas: hash bcrypt.
  \item Usuarios inactivos no pueden autenticarse.
\end{itemize}

\section{Endpoints administrativos}
\begin{tabularx}{\linewidth}{L{3.5cm} L{4.6cm} Y}
  \toprule
  \rowcolor{navy}
  \textcolor{white}{\textbf{Mét.}} &
  \textcolor{white}{\textbf{Ruta}} &
  \textcolor{white}{\textbf{Función}} \\
  \midrule
  GET & \ruta{/admin/dashboard} & Estado del servicio y reindex. \\
  \rowcolor{navysoft}
  GET & \ruta{/admin/branding} & Branding público (sin JWT). \\
  GET/PUT & \ruta{/admin/settings} & LLM, branding, API keys. \\
  \rowcolor{navysoft}
  GET/POST/DEL & \ruta{/admin/documents*} & Biblioteca y uploads. \\
  POST & \ruta{/admin/reindex} & Indexación autenticada. \\
  \rowcolor{navysoft}
  POST & \ruta{/admin/chat/test} & Chat de prueba RAG. \\
  GET/POST/PATCH & \ruta{/admin/users*} & Gestión de administradores. \\
  \bottomrule
\end{tabularx}

\section{Modelo de datos (PostgreSQL)}
Tablas principales:

\begin{itemize}
  \item {admin\_users} — administradores.
  \item {chatbot\_admin\_settings} — fila de configuración (LLM, branding, keys cifradas, flags de reindex).
  \item {admin\_audit\_logs} — auditoría de acciones.
\end{itemize}
% ============================================================
\chapter{Frontend admin}

\section{Stack}
React 18 + TypeScript + Vite. Tipografía Manrope. 

\section{Módulos clave}
\begin{tabularx}{\linewidth}{L{5cm} Y}
  \toprule
  \rowcolor{navy}
  \textcolor{white}{\textbf{Archivo}} &
  \textcolor{white}{\textbf{Rol}} \\
  \midrule
  \ruta{src/api.ts} & Cliente HTTP + JWT en {localStorage}. \\
  \rowcolor{navysoft}
  \ruta{src/auth.tsx} & Contexto de autenticación. \\
  \ruta{src/branding.tsx} & Carga branding y aplica CSS variables. \\
  \rowcolor{navysoft}
  \ruta{src/components/Shell.tsx} & Layout y navegación. \\
  \ruta{src/pages/*.tsx} & Login, Dashboard, Documentos, Settings, Chat, Users. \\
  \bottomrule
\end{tabularx}

\section{Páginas y contratos}
\begin{itemize}
  \item Login consume branding público.
  \item Documentos: upload multipart, delete, reindex con confirmación modal.
  \item Settings: claves (password inputs), modelo, personalización.
  \item Users: no permite desactivar al único admin activo (regla de UI).
\end{itemize}

Variable de entorno del front:
\begin{lstlisting}
VITE_API_BASE_URL=http://127.0.0.1:8000
\end{lstlisting}

% ============================================================
\chapter{Plugin WordPress}

\section{Responsabilidades}
\begin{itemize}
  \item Encolar assets del widget React.
  \item Renderizar el contenedor {\#chatbot-facultad-root}.
  \item Exponer ajustes en el admin de WP (URL API, título, privacidad, enable/disable).
  \item Proxificar el chat hacia la API para no exponer la API directamente al navegador
        en todos los escenarios y aplicar controles propios.
\end{itemize}

\section{Proxy REST}
\begin{itemize}
  \item Namespace: {chatbot-facultad/v1}.
  \item Nonce propio: chatbot\_facultad\_chat.
  \item Límite de longitud de mensaje y rate limit por IP/sesión.
  \item Timeout hacia FastAPI $\sim$90\,s.
  \item Errores hacia el usuario: mensaje genérico amigable.
\end{itemize}

% ============================================================
\chapter{Seguridad}

\section{Controles implementados}
\begin{itemize}
  \item Autenticación JWT en operaciones administrativas.
  \item Contraseñas con bcrypt.
  \item API keys cifradas en reposo (Fernet).
  \item Indexación no pública.
  \item CORS restringido a orígenes configurados.
  \item Auditoría de acciones admin (login, settings, docs, reindex, users, chat test).
  \item Prompt que prohíbe inventar información.
\end{itemize}

% ============================================================
\chapter{Despliegue y operación}

\section{Requisitos}
\begin{itemize}
  \item Python 3.10+, Node 20+, PostgreSQL.
  \item Claves NVIDIA (obligatoria para embeddings) y Gemini si aplica.
\end{itemize}

\section{Pasos resumidos}
\begin{enumerate}
  \item Crear BD {chatbot\_admin} y configurar {.env}.
  \item {pip install -r requirements.txt}
  \item {alembic upgrade head \&\& python -m admin.seed}
  \item {uvicorn api.main:app --host 0.0.0.0 --port 8000}
  \item Frontend: {npm install \&\& npm run build} (o {dev}).
  \item Activar plugin WP y fijar URL base de la API.
\end{enumerate}

\section{Observabilidad}
Logs WARNING/ERROR en {logs/log-YYYY-MM-DD.log}. El endpoint
{GET /health} reporta estado y modelo activo.

% ============================================================
\chapter{Pruebas}

\section{Automatizadas}
\begin{lstlisting}
cd chatbot-facultad
pytest tests/ -v
\end{lstlisting}

Cobertura actual (orientativa):
\begin{itemize}
  \item Auth (login OK/fail/inactive, rutas protegidas).
  \item Settings, branding, cifrado de API keys.
  \item Upload, reindex y chat test (con mocks).
  \item Health y ausencia de {POST /ingest} público.
  \item Factory de proveedores LLM.
\end{itemize}

% ============================================================
\chapter{Contratos de API (resumen)}

\section{Público}
\begin{lstlisting}
GET  /health
POST /chat
Body: { "session_id": "...", "message": "..." }
Resp: { "answer": "...", "sources": [ {file_name, page_label, score?} ] }
\end{lstlisting}

\section{Admin (JWT)}
Documentados en detalle en \ruta{ADMIN.md}. Swagger interactivo:
{/swagger} sobre la API en ejecución.

\begin{figure}[H]
  \centering
  \includegraphics[width=0.92\linewidth]{tec-04-swagger.png}
  \caption{Swagger UI con rutas admin}
  \label{fig:documentos-subida}
\end{figure}
\end{document}
